%%%%%%%%%%%%%%%%%%%%%%%%%%%%%%%%%%%%%%%%%%%%%%%%%%%%%%%%%%%%%%%%%%%%%%
\section{Livrables et contraintes}
%%%%%%%%%%%%%%%%%%%%%%%%%%%%%%%%%%%%%%%%%%%%%%%%%%%%%%%%%%%%%%%%%%%%%%

Cette première phase vise à produire un premier corpus exploitable, et
non une solution définitive de modélisation. Les livrables attendus sont
principalement les suivants : un protocole de collecte, un corpus
partiellement validé, des métadonnées structurées, un pipeline de
transcription et de segmentation, ainsi qu'une base de données prête à
être enrichie dans les phases suivantes.

Le principal enjeu est de rester réaliste dans le cadre du temps
alloué au stage. Il faut donc privilégier une architecture simple,
robuste et évolutive, sans viser une solution trop lourde ou trop
complexe. Les contraintes liées au terrain, à la variabilité dialectale
et à la qualité des enregistrements doivent également être prises en
compte dès le début pour éviter des coûts de correction importants.

%%%%%%%%%%%%%%%%%%%%%%%%%%%%%%%%%%%%%%%%%%%%%%%%%%%%%%%%%%%%%%%%%%%%%%
\section{Points à valider avec l'encadrant}
%%%%%%%%%%%%%%%%%%%%%%%%%%%%%%%%%%%%%%%%%%%%%%%%%%%%%%%%%%%%%%%%%%%%%%

\begin{enumerate}
\item La structure du corpus est-elle suffisamment simple et robuste pour
un premier corpus oral de soussou ?
\item Les métadonnées proposées sont-elles suffisantes pour la
recherche et pour la préparation future des modèles ?
\item Faut-il privilégier Praat, ELAN, MaryTTS ou ESPnet2 pour la
transcription et la segmentation ?
\item L'ordre collecte -- validation -- transcription -- segmentation --
annotation semble-t-il pertinent pour cette première phase ?
\item Les objectifs du stage sont-ils réalistes dans le cadre de la
période de travail prévue ?
\end{enumerate}

\textbf{Conclusion.} Cette note propose une base technique simple,
robuste et directement exploitable pour la constitution d'un premier
corpus oral en soussou. L'objectif principal est de créer une ressource
bien structurée avant toute modélisation. Les choix présentés restent
ouverts à discussion afin d'être ajustés au regard des recommandations
de l'encadrant et des contraintes du projet.